% !TEX program = xelatex
% 컴파일: XeLaTeX 또는 LuaLaTeX (kotex 사용). Overleaf에서 Menu > Compiler를 XeLaTeX로 설정.
% 버전: v1_3 · 2026-08-02
\documentclass[10pt]{article}
\usepackage{kotex}
\usepackage[margin=1in]{geometry}
\usepackage{booktabs}
\usepackage{array}
\usepackage{amsmath}
\usepackage{amssymb}
\usepackage{enumitem}
\usepackage{ragged2e}
\setlist{nosep,leftmargin=1.4em}
\newcolumntype{L}[1]{>{\RaggedRight\arraybackslash}p{#1}}
\renewcommand{\arraystretch}{1.2}

\title{주제 구체화 --- 실시간 진동 결함 진단을 위한\\ 기계 상태·시스템 여유 기반 진단 모드 선택}
\author{이태훈 · v1\_3}
\date{2026-08-02}

\begin{document}
\maketitle

\section*{0. 문서의 목적}
\begin{itemize}
\item 이 문서는 학위논문의 \textbf{주제를 구체화}하기 위한 내부 작업 문서임.
\item 실시간 결함 진단과 elastic scheduling의 조사 결과를 정리하고 연구 공백·정식화·실험 설계·남은 과제를 작성하였음.
\item 적응형 입력(S2)과 실시간 DNN(S4)은 서베이 중이며 현재까지 확보한 범위만 표로 반영하였음.
\end{itemize}

\section{제목 (가제)}
\begin{itemize}
\item \textbf{추천 가제 (국문):} 엣지 디바이스에서 기계 상태와 시스템 여유를 이용한 실시간 진동 결함 진단의 진단 모드 선택
\item \textbf{추천 가제 (영문):} Machine-Condition and Slack-Aware Mode Selection for Real-Time Vibration Fault Diagnosis on Edge Devices
\end{itemize}

\section{Keywords}
\texttt{real-time fault diagnosis}, \texttt{runtime mode selection}, \texttt{elastic scheduling}, \texttt{anomaly-driven adaptation}

\section{국문 초록}
회전 기계의 진동 결함 진단을 엣지 디바이스에서 실시간으로 수행할 때, 진단 신뢰성과 실시간성은 서로 trade-off 관계에 있다. 더 큰 입력 window와 더 무거운 모델은 결함을 더 잘 구분하지만 실행시간과 이용률을 함께 키워 데드라인을 넘길 수 있다. 본 연구는 진단 태스크의 입력 window $W$, 진단 주기를 정하는 hop $H$, 추론 모델 $M$을 하나의 진단 모드로 묶고, 기계 상태 $q$와 시스템 여유 $S$를 함께 기준으로 삼아 매 진단 시점에 모드를 고르는 방법을 제시한다. 정상 상태에서는 가벼운 모드로 여유를 확보하고, 이상 징후가 감지되면 정밀 모드를 고르되 데드라인 안에서 실행 가능한 모드만 후보로 둔다. 이를 통해 검토한 결함 진단 연구가 대체로 평균 latency 중심의 best-effort에 머문 지점을, 데드라인과 실행 가능성 판정과 결합한다.

\section{Abstract}
In real-time vibration fault diagnosis on edge devices, diagnostic reliability and real-time behavior are in a trade-off: a larger input window $W$ and a heavier model $M$ separate faults more reliably but increase execution time and utilization, which can violate the deadline. This work groups the window $W$, the hop $H$ that sets the diagnosis period, and the model $M$ into a single diagnosis mode, and selects one mode at each diagnosis step using both the machine condition $q$ and the system slack $S$. Under a normal condition the runtime chooses a light mode to preserve slack; under a suspicious condition it chooses a precise mode, restricting the candidates to modes that run within the deadline. This couples deadline and feasibility checking to a fault-diagnosis setting where prior work has largely stayed best-effort.

\section{기호 정의}
\begin{center}\footnotesize
\begin{tabular}{@{}lL{11cm}@{}}
\toprule
기호 & 의미 \\ \midrule
$f_s$ & 샘플링 주파수 \\
$W_i,\,H_i,\,M_i$ & $i$번째 후보 모드의 window(샘플), hop(샘플), 모델 \\
$m_i$ & $i$번째 후보 진단 모드, $m_i=(W_i,H_i,M_i)$ \\
$N$ & 후보 모드 개수. 후보 집합 $\{m_1,\dots,m_N\}$ \\
$T_i$ & 진단 주기, $T_i=H_i/f_s$ \\
$C_i$ & 추론 실행시간. 분포로 보고(평균·관측 최대) \\
$I,\,L$ & 간섭 시간, 운영체제 스케줄링 지연 \\
$R_i$ & 관측 응답시간, $R_i=C_i+I+L$ \\
$D$ & 상대 데드라인 \\
$U_i,\,U_{bg},\,U_{bound}$ & 모드 이용률 $U_i=C_i/T_i$, 배경 이용률, 이용률 상한 \\
$q_k$ & 기계 상태 지표. 본 연구에서는 anomaly score \\
$S_k$ & 시스템 여유, $S_k=D-\max(\text{최근 응답시간})$. 평균 대신 관측 최대를 빼서 보수적으로 산정 \\
$Q(m_i,q_k)$ & 현재 기계 상태에서 모드 $m_i$가 달성하는 진단 성능. 진단 정확도(accuracy)로 측정, 필요 시 검출까지 걸린 시간(detection latency)을 함께 봄 \\
$m^{*}_k$ & $k$ 시점에 \textbf{선택된 하나의 모드}. 별표는 선택 결과이며 다른 모드를 포함한다는 뜻이 아님 \\
\bottomrule
\end{tabular}
\end{center}

\section{Background}
\subsection{연구 동기}
\begin{itemize}
\item 회전 기계는 발전설비·생산공정·교통 등의 핵심 요소이므로 결함의 조기 탐지는 안전 확보와 유지보수 비용 절감에 중요하다.
\item 딥러닝 기반 data-driven 진단은 높은 성능을 보이나, 대부분 엣지에서 수집한 데이터를 클라우드로 전송해 추론하는 중앙집중형이라 네트워크 지연·통신 병목으로 실시간 대응이 어렵다. 이에 엣지에서 직접 추론하는 방식이 대안으로 제시된다 [22].
\item 엣지 자원은 제한적이어서 모델 정확도가 높아도 추론 응답시간이 데드라인을 넘거나 지터가 크면 실시간 진단에 쓰기 어렵다. 여기서 진단 신뢰성과 실시간성이 trade-off를 이룬다.
\item KCC 2026 결과가 이 trade-off를 정량적으로 보여준다 [22]. 사용한 FRFconv-TDSNet 모델은 depthwise convolution 계층이 입력 길이에 비례하는 가변 커널 구조를 가져 $O(W^2)$ 복잡도를 갖고, window를 줄이면 연산량이 선형보다 빠르게 감소한다. 데드라인은 한 window의 수집 시간으로 정의된다: $D=W/f_s$. 비중첩 window에서는 진단 주기 $T$가 곧 $D$와 같다.
\end{itemize}
\begin{center}\footnotesize
\begin{tabular}{@{}rrrr@{}}
\toprule
$W$ & 평균 추론 응답시간 & 데드라인 $D=W/f_s$ (8\,kHz) & 이용률 $U=C/D$ \\ \midrule
512 & 40.3\,ms & 64\,ms & 0.63 \\
1024 & 129.8\,ms & 128\,ms & 1.01 \\
2048 & 460.3\,ms & 256\,ms & 1.80 \\
\bottomrule
\end{tabular}
\end{center}
\begin{itemize}
\item 해석: 정밀한 모드로 갈수록 실행시간이 커지고 $T$가 충분히 커지지 않으면 $U>1$이 되어 단일 태스크만으로도 데드라인을 지킬 수 없다. (위 값은 MCU 측정으로 지터가 매우 작아 평균과 최대가 거의 같다. Pi Zero 2W의 Linux에서는 지터가 커질 수 있어 관측 최대를 함께 본다.)
\item 핵심 아이디어: 이 선택을 고정하지 않는다. 정상일 때는 가벼운 모드로 여유를 확보하고, 이상 시 정밀 모드로 전환하되 데드라인 안에서 실행 가능한 경우에만 허용한다. 조절 대상은 window 하나가 아니라 $(W,H,M)$을 묶은 모드다.
\item KCC 모델은 분류 형태이나 내부적으로 anomaly detection을 수행하며 그 anomaly score가 $q$가 된다 [22].
\end{itemize}

\subsection{관련 연구 --- 실시간 결함 진단}
\begin{itemize}
\item 검토한 결함 진단 논문은 데드라인·꼬리 지연·스케줄 가능성을 결합한 사례가 서베이 범위에서 확인되지 않았고 모두 best-effort 수준이었다.
\item Deadline 판정: O(명시+측정), $\triangle$(시간 기준은 있으나 미선언), X(시간 기준 없음).
\end{itemize}
\begin{center}\scriptsize
\begin{tabular}{@{}cL{2.3cm}L{2.4cm}L{1.6cm}L{4.4cm}cL{3.0cm}@{}}
\toprule
\# & 논문 & 문제 상황 & 무엇을 조절 & 무엇으로 정하나 (anomaly detection 방식) & DL & 본 연구와의 차이 \\ \midrule
[1] & Langarica 2020, TASE & 산업 모터 online 진단 & 없음 & DIPCA+SPE 관리한계 초과 시 이상 $\rightarrow$ RBC가 진동 결함 지목 시 CNN (1\,Hz) & X & $q$ 트리거 선례. $S$ 미고려, 모드 선택 없음 \\
[6] & Arciniegas 2025, Discover IoT & 모터 진동 이상 알림 & 없음 & K-means 정상군 거리 임계 초과 시 alert & $\triangle$ & anomaly detection이 alert 조건일 뿐 모드 선택 아님 \\
[2] & Lima 2025, LCIoT & 회전자봉 파손 진단 & 없음 & offline grid search로 결함유형별 최적 $(f_s,W,H)$ 선택 & $\triangle$ & 상태별 최적 $W/H$ 다름 확인. runtime 규칙 없음 \\
[3] & Pubalan 2025, ICSIMA & 베어링 분류(replay) & 없음 & 물리식 $W=60 f_s/\mathrm{RPM}$ & $\triangle$ & 회전속도를 입력 길이로 변환. runtime 적응 아님 \\
[4] & Yang 2023, CSCWD & 베어링 이상 진단(단말-엣지) & 오프로딩 위치 & 단말 AE confidence + 허용 latency 예산으로 엣지 오프로딩 결정 & $\triangle$ & machine evidence+timing 함께 쓰는 최근접 FD. schedulability 없음 \\
[5] & He 2023, TIM & 모터 베어링 음향 진단 & 없음 & 없음(회전속도로 BPFI/BPFO prior만) & $\triangle$ & pipeline 10.294\,s $>$ 1\,s window. ``online $\ne$ 데드라인'' 반례 \\
\bottomrule
\end{tabular}
\end{center}
\begin{itemize}
\item anomaly detection 정량 기준: 통계 관리한계(SPE [1])나 클러스터 거리 임계(K-means [6]). 둘 다 본 연구의 $q$ 임계 설계에 참고가 된다.
\item Lima [2] 표현: 미세 결함이 더 높은 $f_s$·짧은 hop을 요구한 것이 아니라, grid search 결과 그런 설정이 최적으로 선택되었다.
\item Yang [4] 트리거 상세: 단말이 경량 stacked autoencoder를 실행하고, 재구성 신뢰도와 허용 latency 예산을 함께 판단해 무거운 진단을 엣지로 오프로딩할지 결정한다. $q+S$와 형태가 가깝지만 스케줄 가능성 분석은 없다.
\end{itemize}

\subsection{관련 연구 --- Elastic Scheduling}
\begin{itemize}
\item 이론과 적용으로 나눈다. 분류 기준: 이론은 기여가 정형 분석·증명(utilization bound, schedulability, transition 분석)에 있고, 적용은 기여가 실제 플랫폼 구현에 있다. [11], [12]는 제어 응용을 대상으로 하나 기여가 invariant set·weakly-hard 합성 같은 정형 분석이므로 이론으로 분류한다. [13]은 MPSoC·ERIKA RTOS, [14]는 ROS 2 구현이라 적용으로 분류한다.
\item 공통 관찰: 트리거는 시스템 내부 신호나 offline이며 외부 기계 상태 트리거는 확인되지 않았다. 조절 대상은 주기·rate·대역폭으로 정확도를 바꾸지 않는다.
\end{itemize}
\noindent\textbf{이론 논문}
\begin{center}\scriptsize
\begin{tabular}{@{}cL{2.4cm}L{2.6cm}L{2.4cm}L{5.4cm}@{}}
\toprule
\# & 논문 & 조절 대상 & 실시간성 보장 수준 & 핵심 정리 \\ \midrule
[9] & Buttazzo 1998/2002 & task 주기 $T$ & 제한 모델 내 formal hard & elastic task model 원형. task를 spring으로 보고 부하 시 주기를 늘려 이용률을 낮춤 \\
[10] & Salman 2021, ETFA & 앱 주기 + reservation & hard (PRM bound) & 2계층. 앱 수준에서 주기로 먼저 흡수, 실패 시에만 system reservation 조정 \\
[11] & Gifford (Decntr) 2024, RTAS & controller·주기·자원 & hard (invariant set+DBF) & multi-mode 제어를 안전·스케줄 공동설계. 전환 시 carry-over 데드라인 완화 \\
[12] & Xu (Safety-Aware) 2023, RTCSA & 공통 주기, miss 패턴 & soft-but-formal ($\langle m,k\rangle$, $d_{safe}$) & 주기를 공통화하고 안전한 miss 패턴을 automata로 합성 \\
\bottomrule
\end{tabular}
\end{center}
\noindent\textbf{적용 논문 (어떻게 적용했는가)}
\begin{center}\scriptsize
\begin{tabular}{@{}cL{2.2cm}L{2.4cm}L{6.0cm}L{2.4cm}@{}}
\toprule
\# & 논문 & 조절 대상 & 적용 방법 & 실시간성 보장 수준 \\ \midrule
[13] & Burgio 2010, ICCD & TDMA 대역폭+주기 & 중앙 master가 대역폭 요청을 모아 TDMA slot을 재분배, 각 core는 대역폭별 실행시간 LUT를 조회해 elastic 알고리즘으로 주기 재계산(ERIKA RTOS) & soft \\
[14] & Li (ATER) 2025, RTCSA & 센서 샘플링 rate & LTTng로 실행 관측, 메시지 드롭이 늘면 rate$\downarrow$, 실행 여유가 생기면 rate$\uparrow$(ROS 2 Humble, 소스 수정 없음) & empirical only \\
\bottomrule
\end{tabular}
\end{center}

\subsection{관련 연구 --- 적응형 입력·모델, 실시간 DNN}
\begin{itemize}
\item 두 섹션의 구분: 적응형 입력·모델은 무엇을 입력으로 줄지(window·length)를 정하는 데 초점을 두며 분류기가 딥러닝이 아닐 수도 있다(ADW [15]는 SVM). 실시간 DNN 추론은 데드라인 안에서 딥러닝 추론을 어떻게 스케줄·구성할지에 초점을 둔다.
\end{itemize}
\noindent\textbf{적응형 입력·모델}
\begin{center}\scriptsize
\begin{tabular}{@{}cL{2.5cm}L{2.4cm}L{3.8cm}L{1.4cm}L{1.4cm}@{}}
\toprule
\# & 논문 & 조절 대상 & 조절 기준 & 분류기 & 시점 \\ \midrule
[15] & Kim (ADW) 2026, Mathematics & window $W$ ($H=\beta W$) & anomaly deviation score(변동성·주기·국부 스파이크 편차) & SVM & offline \\
[16] & Tang (AIL) 2023, ASC & input length & 물리식(특성 주파수·$f_s$·회전속도) & CNN & offline \\
[17] & Tang (AILWTLN) 2023, EAAI & input length + 경량 모델 & 물리식 $N_a=n^2 f_s/BW$ & 경량 CNN & offline \\
\bottomrule
\end{tabular}
\end{center}
\noindent\textbf{실시간 DNN 추론}
\begin{center}\scriptsize
\begin{tabular}{@{}cL{2.4cm}L{2.6cm}L{3.0cm}L{4.2cm}@{}}
\toprule
\# & 논문 & 조절 대상 & 트리거 & 실시간성 달성 방식 \\ \midrule
[18] & Kang (DNN-SAM) 2022, RTAS & optional subtask 입력 scale & runtime 데드라인 여유 + RoI criticality & 필수/선택 subtask 분리 + non-preemptive EDF 충분조건 + slack reclaim \\
[19] & Li (AMS) 2025, RTCSA & 모델 복잡도·anytime exit & 순간 심박 $HR$ 임계 $\rightarrow D(HR)$ & 상태 기반 모델 선택 + anytime 파라미터 공유 + watchdog \\
[20] & Xu (FLEX) 2024, RTSS & batch·fusion 구성 & 뷰 criticality·GPU 예산·데드라인 & elastic fusion + CEDF 동적 배칭 \\
[21] & Yao (Imprecise) 2020, RTCSA & 실행 depth(필수+선택 단계) & 데드라인 + 선택 단계 성능 기여 & imprecise computation + EDF stage 스케줄 \\
\bottomrule
\end{tabular}
\end{center}

\subsection{연구 공백}
\begin{itemize}
\item 결함 진단 관점: 진동 결함 진단에서 runtime에 $W,H,M$을 함께 조절하고 데드라인·스케줄 가능성을 결합한 사례가 서베이 범위에서 확인되지 않았다.
\item elastic scheduling 관점: elastic scheduling은 정확도를 바꾸는 변수를 다루지 않고 트리거가 시스템 내부나 offline에 한정된다.
\item 두 관점의 결합: $q$와 $S$를 분리된 역할로 사용해 정확도를 바꾸는 실행 가능한 모드를 고르고 스케줄 가능성을 제시한 사례가 서베이 범위에서 확인되지 않았다. 본 연구는 이 결합을 대상으로 한다.
\end{itemize}

\section{Methodology}
\subsection{태스크 모델과 진단 모드}
\begin{itemize}
\item Pi Zero 2W에서 진단 태스크가 주기적으로 실행된다. 매 시점 $k$에 후보 집합 $\{m_1,\dots,m_N\}$에서 모드 하나를 고른다. $m_i=(W_i,H_i,M_i)$, $T_i=H_i/f_s$.
\item $W$와 $H$를 분리하는 이유: ``얼마나 많은 신호 맥락을 보는가($W$)''와 ``얼마나 자주 진단하는가($H$)''는 서로 다른 설계 선택이다.
\item 첫 구현은 $M$ 고정, discrete $(W,H)$만 다룬다. $W$ 후보는 KCC 측정값 $\{512,1024,2048\}$을 기준으로 하고 물리식 하한 $W\ge 60 f_s/\mathrm{RPM}$ [3]을 반영한다.
\end{itemize}
\begin{center}\footnotesize
\begin{tabular}{@{}lrrrrl@{}}
\toprule
모드 & $W$ & $H$ & $T=H/f_s$ & 상대 연산량 ($\propto W^2$) & 주 사용 상황 \\ \midrule
$m_1$ & 512 & 512 & 64\,ms & 1배 (기준) & 정상 \\
$m_2$ & 1024 & 512 & 64\,ms & 4배 & 의심, 진단 빈도 유지 \\
$m_3$ & 1024 & 1024 & 128\,ms & 4배 & 의심, 주기 여유 확보 \\
$m_4$ & 2048 & 2048 & 256\,ms & 16배 & 경고 \\
\bottomrule
\end{tabular}
\end{center}
\begin{itemize}
\item $m_2$와 $m_3$은 같은 window(1024)를 쓰지만 hop이 달라 진단 빈도와 이용률이 다르다. 이것이 $W$와 $H$를 분리하는 이점이다.
\item $N$은 4$\sim$6개를 목표로 하고, offline profiling으로 각 모드 실행시간을 측정해 실행 가능 여부를 미리 표로 만든다.
\end{itemize}

\subsection{실행시간과 이용률}
\begin{itemize}
\item 실행시간 $C_i$는 단일 값이 아니라 분포로 보고한다(평균·관측 최대).
\item PREEMPT\_RT에서 static hard WCET를 증명 없이 주장하지 않고 실측 관측 최대를 기준으로 한다.
\end{itemize}
\begin{equation}
R_i = C_i + I + L,\qquad R_i(\max) \le D,\qquad U_i=\frac{C_i}{T_i},\qquad U_{total}(m_i)=U_{bg}+U_i.
\end{equation}

\subsection{실행 가능 조건과 선택 정책}
\begin{itemize}
\item 첫 연구 질문: 이상 시 정밀 모드로 전환해도 정적으로 스케줄 가능한가, 그 보장을 어떻게 제시하는가.
\item 정책은 실행 가능한 모드 집합 안에서만 진단 성능이 가장 높은 모드를 고른다.
\end{itemize}
\begin{align}
F(k) &= \Big\{\, m_i \;:\; U_{bg}+\tfrac{C_i(\max)}{T_i}\le U_{bound}\ \text{ 그리고 }\ R_i(\max)\le D \,\Big\},\\
m^{*}_k &= \arg\max_{m_i \in F(k)}\; Q(m_i, q_k).
\end{align}
\begin{itemize}
\item $m^{*}_k$는 $k$ 시점에 선택된 하나의 모드이며 다른 모드를 포함한다는 뜻이 아니다.
\item $Q$는 현재 기계 상태에서 모드가 달성하는 진단 성능이며 진단 정확도(accuracy)로 측정한다. 정상 상태에서는 가벼운 모드의 정확도로도 충분하고 이상 상태에서는 정밀 모드의 정확도가 더 높다.
\item $F(k)$가 비면 사전 정의 fallback 모드를 실행한다. 데드라인 miss 후 재실행은 그 자체로 timing 보장이 아니므로, 재실행에 필요한 시간을 미리 확보한 경우에만 둔다.
\item 여유는 평균이 아니라 $S_k=D-\max(\text{최근 응답시간})$으로 정의한다.
\end{itemize}

\subsection{기계 상태와 anomaly detection 기준}
\begin{itemize}
\item anomaly detection은 정상과 비정상을 판별하는 문제이며 결함 유형을 분류하는 fault diagnosis와 구분된다. 본 연구의 $q$는 anomaly detection이 내는 연속 score에서 온다.
\item $q_k$는 anomaly score로 구체화하며 첫 구현은 KCC 모델 내부 anomaly score를 후보로 둔다 [22].
\item 정량 임계 후보: (i) 통계 관리한계 방식(Langarica의 SPE [1]), (ii) 분포 percentile 방식(AMS의 심박 percentile [19]).
\item $q$는 진단 성능 순위를, $S$는 실행 가능한 모드를 정한다. 두 신호의 역할 분리가 정책의 핵심이다.
\end{itemize}

\subsection{알고리즘 (예비 초안, 미확정)}
\noindent\textbf{실행시간을 offline에서 재는 이유}
\begin{itemize}
\item 실행 가능 모드 집합 $F(k)$를 만들려면 지금 실행 중이지 않은 모드의 실행시간도 알아야 한다. 어떤 모드를 실행하지 않고는 그 비용을 관측할 수 없으므로 후보 모드 전체의 실행시간은 offline profiling으로 미리 확보한다.
\item 실제 응답시간은 간섭에 따라 online에서 변하므로 slack $S_k$는 최근 실제 응답시간 이력으로 계산하고 feasibility 판정에도 이 online 값을 반영한다. 즉 offline profile은 초기 기준값, online 최근 이력은 실시간 보정값이다.
\end{itemize}
\noindent\textbf{Offline}
\begin{enumerate}
\item 후보 모드 집합 구성($W$ 후보 + 물리식 하한).
\item 각 모드 실행시간 $C_i$ 측정(평균·최대).
\item 부하 조건별 실행 가능 표 작성.
\item anomaly score 임계(정상·의심·경고)를 검증 데이터로 보정.
\end{enumerate}
\noindent\textbf{Online (매 시점 $k$)}
\begin{enumerate}
\item anomaly score $q_k$ 계산.
\item 최근 응답시간으로 $S_k$ 계산.
\item 실행 가능 모드 집합 $F(k)$ 계산(online 최근 이력 반영).
\item $F(k)$ 안에서 $Q(m_i,q_k)$ 최대 모드 $m^{*}_k$ 선택.
\item $F(k)$가 비면 fallback 실행.
\item 잦은 전환을 막는 hysteresis 적용.
\end{enumerate}

\subsection{$W\cdot H\cdot M$을 함께 묶는 이유}
\begin{center}\footnotesize
\begin{tabular}{@{}cL{4.5cm}L{2.6cm}L{2.6cm}@{}}
\toprule
경우 & 변화 & 이용률 $U$ & 스케줄 가능성 \\ \midrule
1 & $W\downarrow \Rightarrow C\downarrow$, $H,T$ 고정 & 감소 & 개선 \\
2 & $W\downarrow \Rightarrow C\downarrow$, $H,T$도 감소 & 증가 가능 & 악화 가능 \\
\bottomrule
\end{tabular}
\end{center}
\noindent $U_i=C_i/T_i$이므로 window 하나가 아니라 모드 $(W,H,M)$에서 골라야 한다.

\section{Experiment}
\subsection{플랫폼과 검증 실험}
\begin{itemize}
\item 평가 플랫폼: Raspberry Pi Zero 2W(Cortex-A53), Raspberry Pi OS Lite. 추론 backend는 XNNPACK 또는 reference(CMSIS-NN은 Cortex-M 전용).
\item 일반 Linux와 PREEMPT\_RT 비교는 주 실험이 아니라 검증 실험이다. 응답시간·jitter 특성을 파악해 실행시간 모델 신뢰성을 확인한다. cyclictest 분포로 PREEMPT\_RT 실효성을 검증한다.
\end{itemize}

\subsection{baseline과 제안 정책}
\begin{center}\footnotesize
\begin{tabular}{@{}lccL{6.5cm}@{}}
\toprule
정책 & $q$ & $S$ & 설명 \\ \midrule
static-light & X & X & 항상 가벼운 모드 고정 \\
static-precise & X & X & 항상 정밀 모드 고정 \\
q-only & O & X & 기계 상태만 보고 선택, 실행 가능 여부 미확인 \\
S-only & X & O & 여유만 보고 선택, 기계 상태 미반영 \\
제안 ($q+S$) & O & O & 실행 가능한 모드 안에서 진단 정확도가 가장 높은 모드 선택 \\
\bottomrule
\end{tabular}
\end{center}

\subsection{기계 상태 시나리오}
\begin{center}\footnotesize
\begin{tabular}{@{}lll@{}}
\toprule
구간 & anomaly score 후보 임계 & 정책 동작 \\ \midrule
정상 & $<0.3$ & 가벼운 모드로 여유 확보 \\
의심 & $0.3\sim0.7$ & 중간 모드 \\
경고 & $>0.7$ & 정밀 모드, 실행 가능한 경우만 \\
\bottomrule
\end{tabular}
\end{center}
\noindent 임계는 검증 데이터 percentile로 보정한다. 세 구간을 오가는 anomaly score 궤적을 재생해 정책의 모드 열·진단 정확도·데드라인 miss를 baseline과 비교한다.

\subsection{측정 지표}
\begin{itemize}
\item 실시간성: 데드라인 miss 비율, 관측 최대 응답시간, activation jitter, 이용률.
\item 진단: 진단 정확도(accuracy). 필요 시 검출까지 걸린 시간(detection latency).
\item 정책 효과: 제안 정책이 baseline 대비 데드라인 miss를 늘리지 않으면서 진단 정확도를 높이는지, 이상 구간에서 정밀 모드 선택으로 검출 성능이 개선되는지.
\item 오버헤드: 모드 선택은 표 조회·비교라 실행시간 대비 무시 가능할 것으로 예상. ATER [14]처럼 1회 실측해 무시 가능함만 확인하고 지속 측정은 하지 않는다.
\end{itemize}
\noindent\emph{위 값과 결과는 아직 측정하지 않았으며 설계로만 제시한다.}

\section{Novelty}
\begin{itemize}
\item 본 연구의 위치는 두 기준, 트리거와 조절 대상이 만나는 빈 자리다. 외부 기계 상태 트리거와 정확도를 바꾸는 조절 대상 --- 두 자리가 서베이 범위에서 확인되지 않았다.
\end{itemize}
\noindent\textbf{세 가지 차별점}
\begin{enumerate}
\item 조절 대상: 기존은 주기·rate·입력 크기 중 하나 [9--21]. 본 연구는 $W,H,M$ 세 값을 함께 고르며 특히 $W$와 $H$를 분리한다.
\item 트리거: 기존은 시스템 내부 [13, 14] 또는 offline [12]이거나 조건 하나만 본다(slack만 [18], 생리 상태만 [19], anomaly detection만 [1]). 본 연구는 $q$와 $S$를 분리된 역할로 함께 쓴다.
\item 진동 진단 도메인에서의 실시간성 검증: 진동 결함 진단에서 실행 가능한 모드 집합 안에서 선택하고 Pi Zero 2W(Linux, PREEMPT\_RT) 실측 응답시간으로 스케줄 가능성을 확인한다.
\end{enumerate}
\noindent\textbf{가장 가까운 논문 대비}
\begin{center}\footnotesize
\begin{tabular}{@{}L{2.2cm}L{4.0cm}L{6.2cm}@{}}
\toprule
논문 & 무엇을 하나 & 본 연구와 다른 점 \\ \midrule
DNN-SAM [18] & slack으로 입력 scale 선택 & 단일 트리거, 공간 vision·GPU 대상. 본 연구는 시간축 window와 기계 상태를 함께 사용 \\
AMS [19] & 상태로 모델 선택 & 생리 신호 기반, 시스템 여유 미고려. 본 연구는 여유로 실행 가능성을 제한 \\
Langarica [1] & anomaly detection으로 무거운 추론을 gate & 시스템 여유가 없어 모드 선택이 아님. 본 연구는 여유와 결합해 $W,H,M$ 선택 \\
ADW [15] & anomaly 편차로 $W$ 선택 & offline. 본 연구는 runtime에 여유 제약 아래 선택 \\
\bottomrule
\end{tabular}
\end{center}

\section{To-Do}
\begin{itemize}
\item 서베이 보강: 적응형 입력(S2)·실시간 DNN(S4) 조사 완결 후 서베이표와 6.4절 갱신.
\item 정식화 확정: 진단 성능 $Q$ 정의(정확도/검출 시간/$d_{safe}$ [12] 중 적합한 형태), $C(W,M)$ 가변에 맞는 이용률 bound 유도, 모드 집합·$N$·fallback·hysteresis 확정.
\item 실험: Pi Zero 2W 모드 profiling과 feasibility table, Linux vs PREEMPT\_RT 검증, baseline 5종과 제안 정책 비교.
\item 미결: $q$ 정의, 실행 가능 판정 기준(관측 최대/percentile), hysteresis 정도.
\end{itemize}

\begin{thebibliography}{99}\footnotesize
\bibitem{ref1} S. Langarica \emph{et al.}, ``An Industrial Internet Application for Real-Time Fault Diagnosis in Industrial Motors,'' IEEE T-ASE, 2020.
\bibitem{ref2} J. P. B. Lima, ``Real-Time Fault Detection in Induction Motors Using TinyML,'' IEEE LCIoT, 2025.
\bibitem{ref3} B. Pubalan \emph{et al.}, ``Real-Time Bearing Fault Detection and Visualization Using 1D CNN,'' IEEE ICSIMA, 2025.
\bibitem{ref4} C. Yang \emph{et al.}, ``A Novel Bearing Fault Diagnosis Method Based on Stacked Autoencoder and End-Edge Collaboration,'' IEEE CSCWD, 2023.
\bibitem{ref5} C. He \emph{et al.}, ``Real-Time Fault Diagnosis of Motor Bearing via Improved Cyclostationary Analysis onto Edge Computing System,'' IEEE T-IM, 2023.
\bibitem{ref6} S. Arciniegas \emph{et al.}, ``IoT Device for Detecting Abnormal Vibrations in Motors Using TinyML,'' Discover Internet of Things, 2025.
\bibitem{ref7} S. Ma \emph{et al.}, ``A Real-Time Mechanical Fault Diagnosis Approach Based on Lightweight Architecture Search,'' EAAI, 2023.
\bibitem{ref8} 최성현, 김서랑, 김태현, ``저비용 마이크로컨트롤러 환경에서의 경량 딥러닝 기반 회전기계 축 결함 진단 시스템,'' KSC, 2025.
\bibitem{ref9} G. C. Buttazzo \emph{et al.}, ``Elastic Task Model for Adaptive Rate Control,'' IEEE RTSS, 1998; ``Elastic Scheduling for Flexible Workload Management,'' IEEE T-C, 2002.
\bibitem{ref10} S. M. Salman \emph{et al.}, ``Scheduling Elastic Applications in Compositional Real-Time Systems,'' IEEE ETFA, 2021.
\bibitem{ref11} R. Gifford \emph{et al.}, ``Decntr: Optimizing Safety and Schedulability with Multi-Mode Control and Resource Allocation Co-Design,'' IEEE RTAS, 2024.
\bibitem{ref12} S. Xu \emph{et al.}, ``Safety-Aware Implementation of Control Tasks via Scheduling with Period Boosting and Compressing,'' IEEE RTCSA, 2023.
\bibitem{ref13} P. Burgio \emph{et al.}, ``Adaptive TDMA Bus Allocation and Elastic Scheduling,'' IEEE ICCD, 2010.
\bibitem{ref14} R. Li \emph{et al.}, ``ATER: Adaptive Task Execution Rate Regulation for Enhanced Real-Time Performance in ROS 2,'' IEEE RTCSA, 2025.
\bibitem{ref15} M. Kim \emph{et al.}, ``Anomaly Deviation-Based Window Size Selection of Sensor Data for Enhanced Fault Diagnosis Efficiency,'' Mathematics, 2026.
\bibitem{ref16} G. Tang \emph{et al.}, ``Integrating Adaptive Input Length Selection and Unsupervised Transfer Learning for Bearing Fault Diagnosis under Noisy Conditions,'' Applied Soft Computing, 2023.
\bibitem{ref17} G. Tang \emph{et al.}, ``A Novel Transfer Learning Network with Adaptive Input Length Selection and Lightweight Structure,'' EAAI, 2023.
\bibitem{ref18} W. Kang \emph{et al.}, ``DNN-SAM: Split-and-Merge DNN Execution for Real-Time Object Detection,'' IEEE RTAS, 2022.
\bibitem{ref19} Y. Li \emph{et al.}, ``Adaptive Model Selection for Real-Time Heart Disease Detection on Embedded Systems,'' IEEE RTCSA, 2025.
\bibitem{ref20} Y. Xu \emph{et al.}, ``FLEX: Adaptive Task Batch Scheduling with Elastic Fusion in Multi-Modal Multi-View Machine Perception,'' IEEE RTSS, 2024.
\bibitem{ref21} S. Yao \emph{et al.}, ``Scheduling Real-Time Deep Learning Services as Imprecise Computations,'' IEEE RTCSA, 2020.
\bibitem{ref22} 이태훈, 박신형, 김태현, ``Zephyr RTOS 기반 MCU에서의 실시간 기계 결함 진단 추론 성능 분석,'' KCC, 2026. (서울시립대 기계정보공학과)
\end{thebibliography}

\end{document}
